% STATE: DRAFT (awaiting human review).
\clearpage
\section{Main Theorems and Correctness}\label{sec:correctness}

% PROV[user]: Correctness is an umbrella term whose subsystems have different contracts;
% requested a hub-and-spoke diagram, mathematical notation and source references.
% PROV[code]: KILT/LTL/Compile.lean:L62-L83; KILT/Automata/Product.lean:L32-L100;
% KILT/Kripke/Check.lean:L34-L46; KILT/Kripke/Counterexample.lean:L36-L65.
Correctness in KILT is a collection of connected claims.  A compiler must
preserve the meaning of a formula; a product must represent compatible system
and automaton runs; a finite search must decide the existence of an accepting
infinite run.  A returned counterexample needs a further guarantee about its
own execution.  \Cref{fig:correctness-hub} distinguishes these contracts.
The displays below are informal mathematical readings of the named Lean
declarations, rather than new specifications.

% PROV[code]: The four boundaries named in the surrounding paragraph and subsections.
\begin{figure}[ht]
  \centering
  \begin{tikzpicture}
    \node[kiltemphasis,text width=38mm,minimum height=17mm] (hub) at (0,0)
      {KILT correctness\\agreement with\\trace semantics};
    \node[kiltstage,text width=38mm] (language) at (-5.2,0.95)
      {Formula language\\semantics $\leftrightarrow$ acceptance};
    \node[kiltstage,text width=38mm] (graph) at (5.2,0.95)
      {Product and graph\\infinite run $\leftrightarrow$ cycle};
    \node[kiltstage,text width=38mm] (verdict) at (-5.2,-0.95)
      {Checker verdict\\\texttt{true} $\leftrightarrow$ satisfaction};
    \node[kiltstage,text width=38mm] (witness) at (5.2,-0.95)
      {Returned witness\\replayed execution\\violates the property};
    \draw[draw=kiltblue,line width=0.8pt] (hub.north west) -- (language.east);
    \draw[draw=kiltblue,line width=0.8pt] (hub.north east) -- (graph.west);
    \draw[draw=kiltblue,line width=0.8pt] (hub.south west) -- (verdict.east);
    \draw[draw=kiltblue,line width=0.8pt] (hub.south east) -- (witness.west);
  \end{tikzpicture}
  \caption{Different correctness contracts support the public checker and its evidence.}
  \label{fig:correctness-hub}
\end{figure}

\subsection{Semantic and Automata Correspondences}

% PROV[code]: KILT/LTL/NormalForm.lean:L119-L133;
% KILT/LTL/TableauCorrect.lean:L48-L84,L90-L145,L151-L240.
First, correctness of translation means recognising exactly the declared trace
semantics, for every infinite valuation word.  Write $w,i\models\varphi$ for
\lean{Sat} at word $w$ and index $i$, and let $T_f$ be the generalised tableau of an NNF formula
$f$.  \lean{sat_normalize} in \sourcefile{NormalForm.lean} and
\lean{accepts_iff_sat} in \sourcefile{TableauCorrect.lean} establish
\[
  w,i\models\varphi\ \Longleftrightarrow\ w,i\models\operatorname{normalize}(\varphi),
  \qquad
  T_f\text{ accepts }w\ \Longleftrightarrow\ w,0\models f.
\]
The tableau equivalence has two directions.  Soundness shows that a marked
subformula on an accepting run is semantically true.  Its key until lemma,
\lean{until_witness}, obtains a future right-operand witness from recurring
acceptance, excluding indefinite postponement.  Completeness constructs a
state containing exactly the subformulas true at each word position, then
shows that these states form an accepting run.  This semantic state is a
noncomputable proof witness; the executable compiler enumerates finite
assignments instead.

% PROV[code]: KILT/Automata/GeneralizedBuchiCorrect.lean:L65-L135;
% KILT/LTL/Compile.lean:L62-L83.
\sourcefile{GeneralizedBuchiCorrect.lean} supplies the conversion theorem
\lean{toFiniteBuchi_accepts}: adding the round-robin counter neither adds
nor removes accepted words.  The proof lifts an original
run by attaching its counter, or projects a converted run by forgetting it;
counter recurrence connects the two acceptance conditions.  Consequently,
\sourcefile{Compile.lean} provides the model-independent boundary
\[
  B_{\lnot\varphi}\text{ accepts }w
  \quad\Longleftrightarrow\quad \neg(w,0\models\varphi)
  \qquad\text{(\lean{violationAutomaton_accepts}).}
\]
The ordinary B\"uchi automaton is interpreted through CSLib's acceptance
interface.  This establishes the meaning of the violation monitor before any
Kripke structure or graph search is involved.

\clearpage
\subsection{Product and Finite-Search Correctness}

% PROV[user]: Requested a clearer distinction between infinite paths, finite
% state spaces, finite evidence of failure and lasso witnesses.
% PROV[code]: KILT/Kripke/Path.lean, Path, fromState and exists_initial;
% KILT/Kripke/Basic.lean, Serial.
A path has infinitely many positions: $p:\mathbb{N}\to S$.  Finiteness means
that the set $S$ of possible states is finite, not that execution stops after
finitely many steps.  Seriality lets execution continue from every system
state.  A finite route ending at a deadlock is not an infinite path.

% PROV[code]: KILT/Automata/Product.lean:L21-L72;
% KILT/Automata/Kripke.lean, toNA and run_iff_path.
The monitor alone accepts words whether or not the system can produce them.
Product correctness closes this gap.  In \sourcefile{Product.lean},
\lean{product_run_iff} identifies a product execution with an initial system
path together with an automaton run over that path's labels.  The source
label is consumed on each edge, so both sequences refer to the same time
position, including position zero.  With $P=M\otimes A$ and $w_p(i)=L(p(i))$,
\lean{hasAcceptedPath_eq_true_iff} gives
\[
  M.\operatorname{hasAcceptedPath}(A)=\mathtt{true}
  \quad\Longleftrightarrow\quad
  \exists p\in\operatorname{InitialPaths}(M),\ A\text{ accepts }w_p.
\]
% PROV[code]: KILT/Graph/Acceptance.lean:L24-L35,L83-L130;
% KILT/Automata/Product.lean:L54-L72; KILT/LTL/Compile.lean:L80-L83;
% KILT/Automata/ProductLasso.lean, productLassoPath_accepts.
The finite product makes infinite violations detectable by a finite search.
For a graph with starts $I$, edges $\to$ and acceptance $F$,
\sourcefile{Acceptance.lean} proves
\[
  \operatorname{HasAcceptingRun}(I,\to,F)
  \quad\Longleftrightarrow\quad
  \exists q\in F,\ I\leadsto q\ \land\ q\to^{+}q.
\]
If an infinite accepting run exists, finitely many product vertices must
account for infinitely many accepting visits.  Some accepting vertex $q$
therefore occurs at two positions $i<j$.  The route to its first occurrence
gives a prefix; the segment from $i$ to $j$ gives a nonempty return to $q$.
Repeating that segment forever gives an accepting lasso.  Conversely, such a
lasso is an infinite accepting run because each repetition visits $q$.
CSLib's finite recurrence and execution lemmas support this argument.

% PROV[user]: Requested care over why arbitrary infinite counterexamples yield lassos.
% PROV[code]: KILT/Graph/Acceptance.lean, hasAcceptingCycle_iff_hasAcceptingRun;
% KILT/Automata/Product.lean, product_run_iff; KILT/LTL/Compile.lean,
% violationAutomaton_accepts.
Thus, if any violating execution exists, a lasso-shaped violating execution
also exists.  The original execution need not itself settle into a fixed
loop: the proof constructs a repeatable witness.  The repeated vertex is a
system--monitor pair, preserving the monitor's progress as well as the system
state.  Repetition of a system state alone does not establish violation.
Seriality prevents system deadlocks; the lasso argument follows from finite
product states and recurring acceptance.  The positive-length return excludes
an accepting dead end, even when the underlying system is serial.

\subsection{End-to-End Checker Correctness}

% PROV[code]: KILT/Kripke/Check.lean:L16-L46;
% KILT/Automata/Product.lean:L78-L100; KILT/Kripke/Semantics.lean:L17-L24;
% KILT/Kripke/PreparedCheck.lean, check_correct;
% KILT/Graph/Enumeration.lean, runInvariant and runInvariant_eq.
\lean{check_correct} in \sourcefile{KILT/Kripke/Check.lean} composes the
preceding contracts.  Negating the existence of an accepted violation becomes
universal satisfaction of the original formula:
\begin{align*}
  \operatorname{check}(M,\varphi,h_{\mathrm{serial}})=\mathtt{true}
  &\ \Longleftrightarrow\ M\models\varphi\\
  &\ \Longleftrightarrow\
    \forall p\in\operatorname{InitialPaths}(M),\ w_p,0\models\varphi.
\end{align*}
The first direction rules out a false positive; the reverse ensures that a
satisfied formula receives a true verdict.  The companion
\lean{check_eq_false_iff} states that a false verdict is equivalent to the
existence of an initial infinite path whose word violates $\varphi$.
\lean{PreparedCheck.check_correct} establishes the same semantic contract
for the reusable prepared API.  Order invariance makes the Boolean result
independent of the chosen traversal order.

% PROV[code]: KILT/Kripke/Check.lean:L14-L17,L34-L46;
% KILT/Kripke/Finite.lean, structure FiniteKripke;
% KILT/Kripke/Basic.lean, Serial; KILT/Kripke/Semantics.lean:L17-L24.
The public contract assumes finitely many states, decidable equality on states
and atoms, nonempty initials and seriality
($\forall s,\exists t,\ s\to t$).  All initial infinite paths are considered,
without fairness: an available grant need not be taken.  Correctness is
relative to the declared trace semantics; whether a model or formula captures
the intended external system remains a modelling question.

\clearpage
\subsection{Counterexample Guarantees}

% PROV[user]: Requested an explanation of violations settled after finitely many steps.
% PROV[code]: KILT/LTL/Basic.lean, sat_globally and sat_eventually;
% Examples/Models/TwoThreadLock.lean:L116-L140;
% Examples/Models/RequestGrant.lean:L141-L161.
An infinite counterexample can already be decisive after a few steps.
For safety, reaching $\mathit{bad}$ once refutes
$\mathbf{G}\neg\mathit{bad}$ regardless of the remaining execution.
\sourcefile{TwoThreadLock.lean} demonstrates this with a collision at position
six.  For liveness, a finite wait does not refute
$\mathbf{F}\,\mathit{grant}$: the grant may still occur later.
The permanent-waiting path in \sourcefile{RequestGrant.lean} refutes it by
never granting.  A lasso records an infinite continuation in both cases.
The safety failure has a finite witness within its replay, whereas the
waiting loop certifies that granting never occurs.

% PROV[code]: KILT/Graph/Lasso.lean:L15-L38,L65-L70,L94-L128;
% KILT/Automata/ProductLasso.lean:L60-L109;
% KILT/Kripke/Counterexample.lean:L36-L51.
A false verdict proves that a violation exists; extraction must additionally
justify the particular finite object it returns.  A lasso $\ell=u\,v^\omega$
has a finite prefix $u$ and a nonempty loop $v$.  In
\sourcefile{KILT/Kripke/Counterexample.lean},
\lean{counterexample?_valid} establishes the prefix edges, entry and closing
loop edge in the product.  \lean{counterexample?_sound} then concerns the
replay of that same lasso.  Writing $p_\ell$ for its system projection,
its informal contract is
\[
  \operatorname{counterexample?}(M,\varphi,h_{\mathrm{serial}},o)
       =\operatorname{some}(\ell)
  \quad\Longrightarrow\quad
  p_\ell\in\operatorname{InitialPaths}(M)
       \ \land\ \neg(w_{p_\ell},0\models\varphi).
\]
Here $o$ is the system-state enumeration.  Projecting with $\pi_1$ preserves
the exact time sequence, as \lean{productLassoPath_states} in
\sourcefile{ProductLasso.lean} establishes.  The accepting loop entry recurs
forever, so the projected execution violates the formula from position zero.

% PROV[code]: KILT/Kripke/Counterexample.lean:L53-L65;
% KILT/Kripke/PreparedCheck.lean:L95-L120.
Extraction is also complete relative to the Boolean verdict.
\lean{counterexample?_isSome_iff} and
\lean{counterexample?_eq_none_iff}, in the same file, give
\begin{align*}
  \operatorname{counterexample?}(M,\varphi,h_{\mathrm{serial}},o)
       \text{ is some}
    &\ \Longleftrightarrow\
       \operatorname{check}(M,\varphi,h_{\mathrm{serial}})=\mathtt{false},\\
  \operatorname{counterexample?}(M,\varphi,h_{\mathrm{serial}},o)
       =\operatorname{none}
    &\ \Longleftrightarrow\ M\models\varphi.
\end{align*}
Absence therefore establishes satisfaction.  The prepared API has the same
guarantees for its returned lasso; traversal order can change the chosen
witness but preserves its semantic guarantee.

\subsection{Proof Techniques and Trust Boundary}

% PROV[code]: KILT/LTL/NormalForm.lean:L119-L130;
% KILT/LTL/TableauCorrect.lean:L48-L240;
% KILT/Automata/GeneralizedBuchiCorrect.lean:L65-L135;
% KILT/Graph/Acceptance.lean:L83-L130; KILT/Graph/Buchi.lean:L20-L90;
% KILT/Graph/Correspondence.lean, initial_decode, step_decode and reachable_iff;
% KILT/Graph/LassoExtraction.lean, acceptingLasso?_eq_spec.
The proof structure follows these boundaries.  Structural induction handles
normalisation and subformula soundness; temporal arguments recover until
witnesses and propagate release obligations.  Completeness supplies exact
semantic assignments.  Lifting and projection connect the round-robin
representations, while finite recurrence turns infinite acceptance into a
cycle criterion.  For executable refinement,
\lean{hasAcceptingCycle_eq_spec} in \sourcefile{Buchi.lean} connects array
search to the bounded reachability specification.  Dense decoding preserves
starts, edges and reachability, while \sourcefile{LassoExtraction.lean}
relates parent reconstruction to its route specification.  The final checker
proof composes these contracts.

% PROV[code]: Examples/Regression/BuchiTrust.lean:L15-L37 and audit declaration list;
% KILT/Frontend.lean:L14-L46; KILT/LTL/TableauCorrect.lean:L151-L175.
% PROV[build]: 2026-11-15, GUI_LAUNCH=false LEAN_NUM_THREADS=1
% ./scripts/build.sh Examples.Regression.BuchiTrust --verbose, exit 0;
% .lake/logs/build-20261115T090926Z-dzArjQ.log.
% PROV[build]: 2026-11-15, lake env lean --stdin, #print axioms for 13 boundary
% results (normalisation, tableau, conversion, violation language, product/run,
% graph cycle/search, verdict and counterexample contracts), exit 0;
% each reports exactly [propext, Classical.choice, Quot.sound].
Lean checks the resulting proof terms against their declared specifications.
For the main boundary theorems above, direct \texttt{\#print axioms} checks report
exactly \lean{propext}, \lean{Classical.choice} and \lean{Quot.sound}.
\sourcefile{Examples/Regression/BuchiTrust.lean} also checks transitive axiom
dependencies and rejects unfinished proofs or additional axioms; its target
built successfully for this draft.  Classical semantic witnesses support
existence proofs without becoming executable infinite-word tests.  A
\texttt{\#eval} result is execution evidence, whereas \lean{kilt_verify}
constructs a kernel-checked satisfaction proof.  The formal guarantee remains
relative to Lean's kernel, these axioms and the declared definitions; the
examples and rendered counterexamples provide a separate way to inspect
behaviour.
\clearpage
